\scanpage{290}
\mspage{206}
% archive sheet number: 289
qui revient au $=$, \inserted{d'après 4.8.1} si $M$ admet une fega, ou encore une feg canonique.

% slanted
\margin{\uline{\unsure{Commentaires|Commentaire}}}

\deleted{\unsure{4.0}}\unsure{4.8.8}. \deleted{Pour ce qui est de la \unsure{nomenclature}} \inserted{Cette définition est en contradiction avec} \inserted{la terminologie} \deleted{\unsure{suivie}} introduite dans SGA 4 I 9.4 (p 114)\unsure{);|),|;} \deleted{\unsure{q}} \unsure{nous|i.e.} désignons dorénavant sous les noms de ``catégorie \unsure{fbt.|fblt.} $\unsure{\pi|\kappa|k}_0$-accessible'', ``cat. fbt. accessible'' les notions introduites dans 9.4. -- i.e. les ct. admettant une \deleted{\unsure{petite}} \deleted{\unsure{sous-générateurs}} \inserted{sous-catégorie} \unsure{pleines|pleine} \unsure{génératrices|génératrice} (pas $=$ str. génératrice) dont les objets sont $\pi_0$-accessibles (resp. accessibles)

Notons que si $M$ est $\pi_0$-accessible (resp accessible), \inserted{la condition a de} 4.8.1 nous dit que $M$ est \unsure{fb\textsuperscript{=}|fb\textsuperscript{t}|fbt.} $\pi_0$-accessible (resp fbt. accessible), et $=$ qu'il existe une petite sous-catégorie pleine \uuline{strictement} génératrice, qui soit formée d'objets $\pi_0$-accessibles (resp. accessibles). Je doute fort que cette \inserted{dernière} condition \inserted{nécessaire} soit aussi suffisante pour que $M$ soit $\pi_0$-accessible (resp accessible). Mais en vertu de 4.6.3, $3^\circ$ (p 150), il en est ainsi si on suppose $C$ stable par petites $\varinjlim$, ou, ce qui revient au $=$ moyennant $L_{\pi_0}$, \inserted{stable} \unsure{par} limites inductives de cardinal $\leq\pi_0$,
